% GENERATED by scripts/dump_dose_prompts.py -- do not hand-edit.
% Sources of truth: scripts/generate_llm_pairs.py, scripts/generate_meld_pairs.py,
% data/saber/upstream_ref/prompts.py (the prompts that ran).
\section{The prompts, verbatim}
\label{app:prompts}

This appendix prints every prompt behind the paper's attacks word for word, so that a reader can judge how far apart they are rather than take our labels for it. First the four rungs of the ladder of Section~\ref{sec:mechanism}, which change one thing between rungs, the system prompt given to Qwen3-32B-AWQ, while the judge, the source problems and the training settings stay the same; each rung's user message carries only the source problem under a one-line header, and is printed beneath its system prompt. D1--D3 share one output format (D4 asks for restatements only and writes no negatives by design). D1 is MathNet-Retrieve's own published Appendix-F rewriter prompt, changed only so that its output matches our JSON format. Then the two further controls of Appendix~\ref{app:mechanism}, the survey-author prompt D5 and the unrelated near-misses prompt, the judge that filters every LLM-written pair, and the prompts of the MELD and SABER-Math attacks of Section~\ref{sec:ood}. Braces mark the slots the scripts fill.

\paragraph{Word overlap as a distance, and why we do not use it.}
One could measure the distance between prompts by how many words they share. The Jaccard similarity of word sets to D1 is D2 $= 0.411$, D3 $= 0.192$, D4 $= 0.259$, which does not follow the ladder's order: D4 shares more words with D1 than D3 does, yet D4 is the rung whose hard-tier score collapses, because it asks for a different task (restating a problem for another audience) and attaches no negatives. The factorial of Table~\ref{tab:factorial} shows the collapse comes from the missing negatives ($10.30$ hard R@1 over eight seeds once the verified arm's counterexamples are attached, $26.97$ with the recipe arm's own near-misses). Word overlap is therefore the wrong measure of prompt distance, and the hard-tier ordering of Table~\ref{tab:dose-full} does not follow it.

\paragraph{D1: the benchmark's own Appendix-F template (the recipe arm).}
\begin{Verbatim}[breaklines=true,breakanywhere=true,fontsize=\small]
You are a rigorous mathematical editor and rewriter.

## Step 1: Produce {n_pos} Equivalent Variant(s)
- Each is mathematically identical to the original and solvable by the same method.
- Make them look substantially different in structure/phrasing.
- Use at least 3 distinct transformation types across the whole set (may vary by item).
- Keep statements concise.
- For EACH variant, provide:
  > "problem": the LATEX statement,
  > "justification": a 1-sentence explanation of why it's equivalent,
  > "tags": 1-4 short labels naming the transformation(s) used (e.g., "Variable rename", "Modular rewrite").

// Some Transformation Ideas (use multiple or your own):
1. Variable/Parameter Changes (rename, reorder, replace constants with symbols)
2. Algebraic/Arithmetic Rewrites (exponent swap, divisibility <-> congruence, sum <-> product)
3. Language Restatements (synonyms, contrapositive, reorder assumptions)
4. Geometric Equivalences (relabel points, coordinates <-> vectors, area forms)
5. Combinatorial Rephrasings (choose <-> arrange <-> distribute, complement counting)
6. Number Theory Tricks (gcd, modular swaps, prime restatements)
7. Structural Shifts (lemma/theorem framing, reverse flow, regrouping)
8. Sequences <-> Functions (recurrence <-> functional form)
9. Ratios/Normalizations (normalized terms, integer ratios)
10. Duality/Symmetry (polynomial reciprocal, geometry duality, inequality flips)
11. Alternate Representations (closed form <-> recurrence, coordinates <-> trig)
12. Constraint Shifts (additive <-> multiplicative, quantifier swaps)
13. Meta-Level Changes (prove vs. counterexample, classification form)

## Step 2: Produce {n_neg} Near-Miss Variant(s)
- Should look deceptively similar but require a meaningfully different solving method.
- Make small, crucial mathematical changes (e.g., sign flip, add a square, exponent shift, modulus tweak, relation type change, altered bounds).
- Ensure each truly alters the strategy (not a trivial edit).
- For EACH near-miss, provide:
  > "problem": the LATEX statement,
  > "justification": a 1-sentence explanation of why it's a near miss,
  > "tags": 1-4 short labels naming the modification(s) (e.g., "Sign flip", "Modulus change", "Inequality direction").

## Step 3: Return Strict JSON
Return only a single valid JSON object using this exact schema:
{{
  "original_problem": "LaTeX string of the cleaned original problem",
  "equivalent_variants": [
    {{"problem": "...", "justification": "...", "tags": ["...", "..."]}}
  ],
  "near_miss_variants": [
    {{"problem": "...", "justification": "...", "tags": ["..."]}}
  ]
}}
// Rules:
- All math must be in LATEX.
- Exactly {n_pos} equivalent and {n_neg} near-miss variants.
- Each variant MUST include non-empty "tags".
- Output JSON only; no code fences, no extra text.
- If your draft is not valid JSON, fix it until it is valid and return only the JSON.
\end{Verbatim}

User message of D1:
\begin{Verbatim}[breaklines=true,breakanywhere=true,fontsize=\small]
Original problem:

{problem}
\end{Verbatim}

\paragraph{D2: a close paraphrase of that template.}
\begin{Verbatim}[breaklines=true,breakanywhere=true,fontsize=\small]
You act as a careful, exacting editor of mathematics problems.

Your assignment has two parts.

### Part A -- equivalent restatement(s)
Write {n_pos} variant(s) of the problem you are given that carry exactly the same mathematical content as the original: identical in substance and solvable by the very same method, yet visibly different in wording and surface structure. Across the whole set of variants, draw on at least 3 different kinds of transformation; options include renaming variables, reordering them, or swapping constants for symbols; rewriting the algebra or arithmetic (exchanging exponents, trading divisibility statements for congruences, turning sums into products); restating the language with synonyms, a contrapositive, or reordered assumptions; relabeling geometric points or moving between coordinates, vectors, and area formulas; recasting counting arguments (choosing vs. arranging vs. distributing, counting the complement); number-theoretic rewrites via gcd, moduli, or primes; reframing as a lemma or reversing the logical flow; trading recurrences for functional forms; normalizing ratios; exploiting duality or symmetry; switching to an alternate representation; shifting constraints between additive and multiplicative form or swapping quantifiers; or changing the meta-level (prove vs. find a counterexample, classification form). You may also invent transformations of your own. Keep every statement concise. Record for each variant:
  > "problem": the statement in LATEX,
  > "justification": a single sentence explaining why nothing mathematical changed,
  > "tags": 1-4 brief labels naming the transformation(s) applied (e.g., "Variable rename", "Modular rewrite").

### Part B -- near-miss lookalike(s)
Write {n_neg} problem(s) that closely resemble the original at a glance yet genuinely demand a different solving method. Achieve this with small but decisive mathematical alterations -- a flipped sign, an added square, a shifted exponent, an adjusted modulus, a changed relation type, or moved bounds -- and make certain each alteration truly forces a new strategy rather than being a cosmetic edit. Record for each near-miss:
  > "problem": the statement in LATEX,
  > "justification": a single sentence explaining why it merely looks similar,
  > "tags": 1-4 brief labels naming the modification(s) made (e.g., "Sign flip", "Modulus change", "Inequality direction").

### Part C -- answer in strict JSON
Return only a single valid JSON object using this exact schema:
{{
  "original_problem": "LaTeX string of the cleaned original problem",
  "equivalent_variants": [
    {{"problem": "...", "justification": "...", "tags": ["...", "..."]}}
  ],
  "near_miss_variants": [
    {{"problem": "...", "justification": "...", "tags": ["..."]}}
  ]
}}
Constraints: express all mathematics in LATEX; produce exactly {n_pos} equivalent and {n_neg} near-miss variant(s); never leave "tags" empty; emit nothing besides the JSON object -- no code fences, no commentary; and if your draft fails to parse as JSON, repair it until it is valid and return only the JSON.
\end{Verbatim}

User message of D2:
\begin{Verbatim}[breaklines=true,breakanywhere=true,fontsize=\small]
The problem to work from:

{problem}
\end{Verbatim}

\paragraph{D3: a rewriter prompt in a deliberately different style.}
\begin{Verbatim}[breaklines=true,breakanywhere=true,fontsize=\small]
Hi! I coach a competition-math training group, and I'm assembling this week's worksheet. I need your help with one exercise at a time.

Here's the game. My students have already seen the problem I'm about to show you. To test whether they recognize a problem by its mathematics rather than by its looks, the worksheet mixes disguises of it in among traps:

1. {n_pos} disguise(s): the same problem dressed up to look new. A disguise must be the original in every mathematical respect -- same content, same answer, same solution path -- but a returning student shouldn't recognize it at first sight. Dress it up however you like: new letters, a different setting or story, reshuffled givens, reworded conditions. Keep it lean; competition problems don't ramble.

2. {n_neg} trap(s): problems that LOOK like today's problem but are NOT it. A good trap differs by one quiet mathematical detail -- perhaps a sign, a power, a modulus, a bound, or the type of relation -- chosen so that the trap sends a solver down a genuinely different path. A trap that ends up solved the same way as the original is a failed trap; don't hand me those.

For every disguise and every trap, jot down:
  > "problem": the full statement (all mathematics in LATEX),
  > "justification": one line -- for a disguise, why it's still the same problem; for a trap, why it only looks like it,
  > "tags": 1-4 short notes on what you changed.

So my worksheet tool can ingest your work, reply with strict JSON only:
{{
  "original_problem": "LaTeX string of the cleaned original problem",
  "equivalent_variants": [
    {{"problem": "...", "justification": "...", "tags": ["...", "..."]}}
  ],
  "near_miss_variants": [
    {{"problem": "...", "justification": "...", "tags": ["..."]}}
  ]
}}
House rules: disguises go under "equivalent_variants" and traps under "near_miss_variants", exactly {n_pos} and {n_neg} of them; every "tags" list non-empty; all math in LATEX; nothing outside the JSON object -- no code fences, no chat; and if what you wrote isn't valid JSON, fix it before you answer. Thanks!
\end{Verbatim}

User message of D3:
\begin{Verbatim}[breaklines=true,breakanywhere=true,fontsize=\small]
Today's problem for the worksheet:

{problem}
\end{Verbatim}

\paragraph{D4: a restatement prompt unrelated to the recipe.}
\begin{Verbatim}[breaklines=true,breakanywhere=true,fontsize=\small]
You are an editor adapting mathematics problems for textbooks aimed at different audiences.

Rewrite the given problem for a student at a DIFFERENT level of mathematical maturity than its original audience. Pick whichever direction suits the problem: a younger student meeting the topic for the first time (friendlier vocabulary, gentler sentence rhythm, concrete framing), or an advanced undergraduate who prefers terse, formal statements. The restated problem must remain mathematically identical to the original -- the same question, the same given conditions, the same answer, solvable by the same reasoning. You may rename characters or objects in story contexts and adapt notation conventions to the audience, but you must not add, remove, weaken, or strengthen any mathematical condition.

Produce {n_pos} restatement(s). For each, record:
  > "problem": the restated problem (all mathematics in LATEX),
  > "audience": who you wrote it for (e.g., "middle-school student", "advanced undergraduate"),
  > "justification": one sentence confirming the mathematics is unchanged.

Return only a single valid JSON object using this exact schema:
{{
  "original_problem": "LaTeX string of the cleaned original problem",
  "restatements": [
    {{"problem": "...", "audience": "...", "justification": "..."}}
  ]
}}
// Rules:
- All math must be in LATEX.
- Exactly {n_pos} restatement(s).
- Output JSON only; no code fences, no extra text.
- If your draft is not valid JSON, fix it until it is valid and return only the JSON.
\end{Verbatim}

User message of D4:
\begin{Verbatim}[breaklines=true,breakanywhere=true,fontsize=\small]
Problem to adapt:

{problem}
\end{Verbatim}

\paragraph{D5: the survey-author prompt, the second recipe-free reference.}
A restatement prompt whose task framing is not an audience change: the model writes the problem as a survey author would present another author's problem. Like D4 it asks for restatements only.
\begin{Verbatim}[breaklines=true,breakanywhere=true,fontsize=\small]
You are a mathematician writing the problems section of a survey article that collects competition problems on one theme.

Restate the given problem in your own words as a self-contained statement, the way you would present another author's problem inside your survey: your own sentence structure, your own notation and your own order of presentation, with the mathematics unchanged. Do not simplify or generalise the problem, do not add hints, and do not change its level.

Produce {n_pos} restatement(s). For each, record:
  > "problem": the restated problem (all mathematics in LATEX),
  > "context": one short phrase naming the theme under which your survey files this problem (e.g., "bounds on sums of reciprocals"),
  > "justification": one sentence confirming the mathematics is unchanged.

Return only a single valid JSON object using this exact schema:
{{
  "original_problem": "LaTeX string of the cleaned original problem",
  "restatements": [
    {{"problem": "...", "context": "...", "justification": "..."}}
  ]
}}
// Rules:
- All math must be in LATEX.
- Exactly {n_pos} restatement(s).
- Output JSON only; no code fences, no extra text.
- If your draft is not valid JSON, fix it until it is valid and return only the JSON.
\end{Verbatim}

User message of D5:
\begin{Verbatim}[breaklines=true,breakanywhere=true,fontsize=\small]
Problem to restate for the survey:

{problem}
\end{Verbatim}

\paragraph{The unrelated near-misses prompt.}
D4's restatement task with a second step that writes minimal-edit companions, so that near-misses can be produced by a prompt that shares nothing with the recipe.
\begin{Verbatim}[breaklines=true,breakanywhere=true,fontsize=\small]
You are an editor adapting mathematics problems for textbooks aimed at different audiences.

First, rewrite the given problem for a student at a DIFFERENT level of mathematical maturity than its original audience. Pick whichever direction suits the problem: a younger student meeting the topic for the first time (friendlier vocabulary, gentler sentence rhythm, concrete framing), or an advanced undergraduate who prefers terse, formal statements. The restated problem must remain mathematically identical to the original -- the same question, the same given conditions, the same answer, solvable by the same reasoning. You may rename characters or objects in story contexts and adapt notation conventions to the audience, but you must not add, remove, weaken, or strengthen any mathematical condition.

Second, this textbook pairs every exercise with a short "spot the difference" drill. Write {n_neg} companion exercise(s) for that drill. Each companion should read like the original problem at a glance, in the same register and about the same length, yet ask something mathematically different, so that a student who solves it by copying the original's reasoning arrives at a wrong answer. Change one thing that matters (a given condition, a quantity, a relation between the objects, or what is asked for) and leave everything else as it was.

Produce {n_pos} restatement(s) and {n_neg} companion(s). For each restatement, record:
  > "problem": the restated problem (all mathematics in LATEX),
  > "audience": who you wrote it for (e.g., "middle-school student", "advanced undergraduate"),
  > "justification": one sentence confirming the mathematics is unchanged.
For each companion, record:
  > "problem": the companion exercise (all mathematics in LATEX),
  > "what_changed": one sentence naming the change and why the answer differs.

Return only a single valid JSON object using this exact schema:
{{
  "original_problem": "LaTeX string of the cleaned original problem",
  "restatements": [
    {{"problem": "...", "audience": "...", "justification": "..."}}
  ],
  "companions": [
    {{"problem": "...", "what_changed": "..."}}
  ]
}}
// Rules:
- All math must be in LATEX.
- Exactly {n_pos} restatement(s) and {n_neg} companion(s).
- Output JSON only; no code fences, no extra text.
- If your draft is not valid JSON, fix it until it is valid and return only the JSON.
\end{Verbatim}

User message:
\begin{Verbatim}[breaklines=true,breakanywhere=true,fontsize=\small]
Problem to adapt:

{problem}
\end{Verbatim}

\paragraph{The judge that keeps or drops every LLM-written pair.}
Ours, since the benchmark does not print its judge prompt; run through Qwen3-32B-AWQ, and for the two-judge control of Appendix~\ref{app:mechanism} through Mistral-Small-24B-Instruct-2501 as well.
\begin{Verbatim}[breaklines=true,breakanywhere=true,fontsize=\small]
You are an expert mathematical olympiad editor acting as a strict judge of mathematical equivalence.

You will be given Problem A and Problem B. Decide whether they are MATHEMATICALLY EQUIVALENT in the strict sense (Invariance): B is a reformulation of A such that the two problems have identical mathematical content — every complete solution of one converts to a complete solution of the other by mechanical translation (renaming variables, rewriting notation, restating language), and both are solved by the same method. Superficial similarity is NOT equivalence: a sign flip, changed exponent, altered modulus, different bound, or added/removed constraint that changes the answer or the solving strategy makes them NOT equivalent.

Think through the mathematics carefully, then return ONLY a single valid JSON object:
{"verdict": "equivalent" | "not_equivalent", "confidence": 0.0-1.0, "reason": "one-sentence justification"}
// Rules:
- Output JSON only; no code fences, no extra text.
- "confidence" is your subjective probability that your verdict is correct.
\end{Verbatim}

User message of the judge:
\begin{Verbatim}[breaklines=true,breakanywhere=true,fontsize=\small]
Problem A:

{anchor}

Problem B:

{candidate}
\end{Verbatim}

\paragraph{The MELD attack: the reconstructed generation prompt.}
MELD publishes the prompts of its training pipeline but not the one behind its evaluation set; this is our reconstruction of the procedure its paper describes, into the schema of its released records. The connection descriptions per subfield pair are ours.
\begin{Verbatim}[breaklines=true,breakanywhere=true,fontsize=\small]
You are a research mathematician building an evaluation set of mathematically equivalent but lexically different statement pairs.

You are given two complementary mathematical subfields and a description of the connection between them. Produce exactly {n_pairs} statement pairs. In every pair:
- entry_1 is written entirely in the language of {field_1};
- entry_2 is written entirely in the language of {field_2};
- the two statements are mathematically equivalent: they assert the same thing about the same underlying mathematics, and a reader who knows both dialects would call them the same statement;
- they cannot be matched by surface-level lexical similarity: share as few words, symbols and notational conventions as possible while staying faithful;
- the content belongs to a basic graduate curriculum in mathematics.

Each pair carries a short "topic" label naming the shared idea. All {n_pairs} topics must be distinct from each other.

Write each statement in natural language and LaTeX, in the voice of a textbook definition or theorem statement, one or two sentences long.

Output ONLY a single valid JSON object, no code fences and no commentary:
{{"pairs": [{{"topic": "<short label>", "entry_1": {{"framing": "{field_1}", "statement": "<statement>"}}, "entry_2": {{"framing": "{field_2}", "statement": "<statement>"}}}}, ...]}}
\end{Verbatim}

User message (first call per subfield pair):
\begin{Verbatim}[breaklines=true,breakanywhere=true,fontsize=\small]
Subfield 1: {field_1}
Subfield 2: {field_2}

Connection between the two fields: {connection}

Produce {n_pairs} pairs.
\end{Verbatim}

User message (later calls, which pass back the topics already produced):
\begin{Verbatim}[breaklines=true,breakanywhere=true,fontsize=\small]
Subfield 1: {field_1}
Subfield 2: {field_2}

Connection between the two fields: {connection}

The following topics have already been written for this subfield pair:
{covered}

Produce {n_pairs} further pairs. Prefer topics that are not in that list; where the shared idea is unavoidable, write a genuinely different statement of it (different objects, different hypotheses, different phrasing).
\end{Verbatim}

\paragraph{The MELD attack: the rewritten prompt of the $84/16$ control.}
\begin{Verbatim}[breaklines=true,breakanywhere=true,fontsize=\small]
You are preparing translation exercises for a graduate reading course that runs two seminars in parallel.

Students in seminar A work in {field_1}; students in seminar B work in {field_2}. Write exactly {n_pairs} exercises. Each exercise states one fact twice, once as seminar A would record it and once as seminar B would record it, so that a student attending only one seminar still meets the same mathematics.

For each exercise:
- entry_1 is the seminar-A version, phrased as that seminar phrases things;
- entry_2 is the seminar-B version, phrased as that seminar phrases things;
- both versions must be true and must say the same thing;
- the content should sit within a basic graduate curriculum.

Give each exercise a short "topic" label naming the fact being recorded, and make the {n_pairs} labels distinct.

Output ONLY a single valid JSON object, no code fences and no commentary:
{{"pairs": [{{"topic": "<short label>", "entry_1": {{"framing": "{field_1}", "statement": "<statement>"}}, "entry_2": {{"framing": "{field_2}", "statement": "<statement>"}}}}, ...]}}
\end{Verbatim}

\paragraph{The MELD attack: the validity and equivalence screen.}
MELD's three-question check, run through Qwen3-32B-AWQ; we gate on validity and equivalence and only record the dissimilarity answer.
\begin{Verbatim}[breaklines=true,breakanywhere=true,fontsize=\small]
You are an expert mathematician screening candidate items for an evaluation set of mathematically equivalent but lexically different statement pairs.

You will be given Statement 1 (written in one subfield's language) and Statement 2 (written in another's). Answer three questions:
(i)   is each statement, on its own, a valid and correct mathematical statement?
(ii)  are the two statements mathematically equivalent -- do they assert the same thing about the same underlying mathematics?
(iii) could the pair be rewritten to sound LESS similar to each other without breaking the equivalence?

Judge (ii) strictly. Sharing a topic is not equivalence: a changed hypothesis, a dropped finiteness assumption, a different quantifier, or a statement that merely implies the other in one direction makes them NOT equivalent.

Return ONLY a single valid JSON object, no code fences and no commentary:
{"entry_1_valid": true|false, "entry_2_valid": true|false, "equivalent": true|false, "could_be_less_similar": true|false, "confidence": 0.0-1.0, "reason": "one sentence"}
\end{Verbatim}

User message of the screen:
\begin{Verbatim}[breaklines=true,breakanywhere=true,fontsize=\small]
Statement 1 ({framing_1}):

{statement_1}

Statement 2 ({framing_2}):

{statement_2}
\end{Verbatim}

\paragraph{The SABER-Math attack: SABER's own solution-summary prompt.}
Taken byte for byte from SABER's repository (\texttt{annotate/ideas/prompts.py}) and run through Qwen3-32B-AWQ in place of GPT-OSS-120B; the summaries it writes are the only text the attack trains on.
\begin{Verbatim}[breaklines=true,breakanywhere=true,fontsize=\small]
# INSTRUCTION
You are an expert mathematical anotator tasked with identifying the *core idea* - the central mathematical insight, from a math problem.

# GOALS
1. Identify what makes the solution work conceptually, not how to carry it out. Capture the untrial step or idea that is the greatest hint for the solution.
2. Never include any multi-step reasoning, equations or numeric computations. Don't include any anotations that are in the solution, but not in the original problem statement.
3. Never try to solve the problem on your own and don't include your reasoning or thoughs.
4. Output a sinlge valid JSON object matching the schema below.
5. Structure the ideas imperatively so they look like you are giving a hint to someone.
5. If the problem seems too easy, straightforward or you can't identify a core idea, store its value as 'null' and set the 'noCoreIdea' to 'true'

# SCHEMA
```json
{
    "noCoreIdea": <true|false>,
    "coreIdea": "<string - one short sentence (up to 30 words) naming the main insight to the problem>",
    "supporingIdeas: ["<strings - 0-3 short techniques phrases>"],
    "keywords": ["<strings - 1-2 word phrases summarizing the ideas, theorems, etc. in the solution"],
    "confidence": <0.0-1.0>
}

Here are the problem statement and solution:

Statement: {{statement}}

Solution: {{solution}}
\end{Verbatim}

